\section{A bounded integration protocol}\label{sec:C21:integration}
The integration step should be designed as carefully as worker execution. A manager receives candidate artifacts, validates their provenance, identifies disagreements, and chooses whether more work is justified. It should not simply concatenate worker summaries and ask another model to produce a fluent answer. Concatenation preserves text but does not resolve incompatible assumptions, duplicated evidence, or different interpretations of the task.

One useful protocol maintains a claim graph. Nodes are claims or artifact predicates; edges identify supporting evidence, contradiction, or dependency. Workers submit proposed graph updates rather than final conclusions alone. The manager checks whether two claims concern the same population, version, and resource envelope before labeling them contradictory. Apparent disagreement can disappear when one worker reports author results and another reports an independent reproduction under a different environment. Genuine unresolved disagreement should remain visible.

Integration can stop when all required predicates have accepted evidence, unresolved contradictions have been represented, and the remaining expected value of additional work is below its cost under the declared objective. This stopping rule is conceptual because value estimates are uncertain. An implementation can approximate it with bounded rounds, explicit missing-evidence categories, and a limit on repeated requests for the same information. The approximation should be evaluated rather than presented as optimal metareasoning.

The protocol also needs a quarantine path. A worker result containing unverifiable citations, unexplained external changes, or suspicious instructions should be retained for diagnosis without entering trusted shared memory. Quarantine is not deletion and does not itself prove maliciousness. It prevents uncertainty about one artifact from silently becoming authority for subsequent actions. False quarantines reduce useful throughput, so their rate belongs in the evaluation alongside successful containment.

For a scientific review task, integration should retain disagreements over contamination, baseline strength, and statistical uncertainty even when the final narrative would be simpler without them. For a coding task, it should retain a failing regression test rather than allow a worker's confident explanation to supersede it. For enterprise actions, it should recheck authorization at the effect boundary. These examples share a principle: evidence and authority are properties of records and current policy, not consequences of a worker's role label.

The resulting team is best understood as a controlled distribution of computation and evidence. Parallelism supplies opportunities for coverage and latency reduction. Typed handoffs preserve what each opportunity discovered. Shared-state protocols prevent those discoveries from becoming incoherent or unauthorized. Independent evaluation determines whether the arrangement improves useful outcomes. None of these layers can be omitted by appealing to the number of agents or the apparent sophistication of their conversation.
